\subsection{Aplicación a los morfismos propios: III. Invariancia de la
característica de Euler-Poincaré y del polinomio de Hilbert}
\label{subsection:III.7.9}
\label{III.7.9}

\begin{env}[7.9.1]
\label{III.7.9.1}
Sean $A$ un anillo y $M$ un $A$-módulo proyectivo de tipo finito.
Recordemos (Bourbaki, \emph{Alg.~comm.}, cap.~II, \S5, n.\textsuperscript{o}~2)
que esto equivale a decir que el $\mathscr O_X$-módulo $\widetilde M$
asociado a $M$ sobre $X=\Spec(A)$ es localmente libre de tipo finito.
Para todo $\mathfrak p\in\Spec(A)$, se llama \emph{rango de $M$ en
$\mathfrak p$}, y se denota por $\operatorname{rank}_{\mathfrak p}(M)$,
al rango del $A_{\mathfrak p}$-módulo libre $M_{\mathfrak p}$ (o,
equivalentemente, al rango en $\mathfrak p$ del $\mathscr O_X$-módulo
localmente libre $\widetilde M$). Por tanto,
\[
  \operatorname{rank}_{\mathfrak p}M
  =
  \operatorname{rank}_{\mathfrak p}(M_{\mathfrak p})
  =
  \operatorname{rank}_{\kres(\mathfrak p)}
    \bigl(M\otimes_A\kres(\mathfrak p)\bigr).
  \tag{7.9.1.1}\label{III.7.9.1.1}
\]
\end{env}

\begin{proposition}[7.9.2]
\label{III.7.9.2}
Sea $P_\bullet$ un complejo finito de $A$-módulos proyectivos de
tipo finito y, para todo $A$-módulo $M$, pongamos
\[
  T_\bullet(M)=H_\bullet(P_\bullet\otimes_A M).
\]
Entonces, para todo $\mathfrak p\in\Spec(A)$, se tiene
\[
  \sum_i(-1)^i
    \operatorname{rank}_{\kres(\mathfrak p)}
      T_i\bigl(\kres(\mathfrak p)\bigr)
  =
  \sum_i(-1)^i\operatorname{rank}_{\mathfrak p}(P_i).
  \tag{7.9.2.1}\label{III.7.9.2.1}
\]
\end{proposition}

\begin{proof}
En efecto, por definición,
$T_i(\kres(\mathfrak p))
 =H_i(P_\bullet\otimes_A\kres(\mathfrak p))$.
Teniendo en cuenta \eqref{III.7.9.1.1}, la fórmula
\eqref{III.7.9.2.1} no es sino la invariancia de la característica
de Euler-Poincaré de un complejo finito de espacios vectoriales de
dimensión finita al pasar a la homología
\sref[0]{0.11.10.2}.
\end{proof}

\begin{corollary}[7.9.3]
\label{III.7.9.3}
La función
\[
  \mathfrak p\longmapsto
  \sum_i(-1)^i
    \operatorname{rank}_{\kres(\mathfrak p)}
      T_i\bigl(\kres(\mathfrak p)\bigr)
\]
es localmente constante en $\Spec(A)$.
\end{corollary}

\begin{theorem}[7.9.4]
\label{III.7.9.4}
Sean $Y$ un preesquema localmente noetheriano,
$f:X\to Y$ un morfismo propio y $\mathscr P_\bullet$ un complejo
finito de $\mathscr O_X$-módulos coherentes y planos sobre $Y$.
Si se pone
\[
  \mathscr T_\bullet(\mathscr M)
  =
  \mathscr H^{-\bullet}
    \bigl(f,\mathscr P^\bullet
      \otimes_{\mathscr O_Y}\mathscr M\bigr)
  \qquad\text{(cf.~\eqref{III.7.7.1.1})},
\]
entonces la función

\oldpage[III]{209}
\[
  y\longmapsto
  \sum_i(-1)^i
    \operatorname{rank}_{\kres(y)}
      \mathscr T_i\bigl(\kres(y)\bigr)
  \tag{7.9.4.1}\label{III.7.9.4.1}
\]
es localmente constante en $Y$.
\end{theorem}

\begin{proof}
Podemos limitarnos al caso en que $Y=\Spec(A)$ es afín y $A$ es
noetheriano. Como el complejo $\mathscr P_\bullet$ es finito,
\sref{III.7.7.12}{(i)} muestra que
\[
  \mathscr T_p(\mathscr M)
  =
  \mathscr H_p(\mathscr L_\bullet
    \otimes_{\mathscr O_Y}\mathscr M),
\]
donde $\mathscr L_\bullet=\widetilde{L}_\bullet$ y
$L_\bullet$ es un complejo finito de $A$-módulos proyectivos de
tipo finito. El teorema se deduce entonces de
\sref{III.7.9.3}.
\end{proof}

\begin{env}[7.9.5]
\label{III.7.9.5}
Bajo las hipótesis de \sref{III.7.9.4}, la función
\eqref{III.7.9.4.1} es \emph{constante} cuando $Y$ es
\emph{conexo}. Cuando $Y$ es conexo y no vacío, su único valor
(entero) se denota por
\[
  \operatorname{EP}(f,\mathscr P_\bullet),\qquad
  \operatorname{EP}(Y,\mathscr P_\bullet),
  \qquad\text{o simplemente}\qquad
  \operatorname{EP}(\mathscr P_\bullet)
\]
cuando no haya riesgo de confusión, y se dice que este entero es la
\emph{característica de Euler-Poincaré de $\mathscr P_\bullet$}
relativa a $f$ (o a $Y$). En general, se escribe asimismo
\[
  \operatorname{EP}(f,\mathscr P_\bullet;y),\qquad
  \operatorname{EP}(Y,\mathscr P_\bullet;y),\qquad
  \operatorname{EP}(\mathscr P_\bullet;y)
\]
para el segundo miembro de \eqref{III.7.9.4.1}.
\end{env}

\begin{env}[7.9.6]
\label{III.7.9.6}
Bajo las hipótesis de \sref{III.7.9.4} relativas a
$X$, $Y$ y $f$, sea
\[
  0\longrightarrow\mathscr P'_\bullet
    \xrightarrow{u}\mathscr P_\bullet
    \xrightarrow{v}\mathscr P''_\bullet
    \longrightarrow 0
\]
una sucesión exacta de complejos finitos de $\mathscr O_X$-módulos
coherentes y planos sobre $Y$, donde los homomorfismos $u$ y $v$
son de grados pares $2d$ y $2d'$, respectivamente.
Como $\mathscr T_\bullet$ es un funtor homológico
\sref{III.7.7.1}, se tiene una sucesión exacta de homología
\[
  \cdots\longrightarrow
  \mathscr T_i(\mathscr P'_\bullet,\kres(y))
  \longrightarrow
  \mathscr T_{i+2d}(\mathscr P_\bullet,\kres(y))
  \longrightarrow
  \mathscr T_{i+2d+2d'}(\mathscr P''_\bullet,\kres(y))
  \longrightarrow
  \mathscr T_{i-1}(\mathscr P'_\bullet,\kres(y))
  \longrightarrow\cdots ,
\]
que, además, solo tiene un número finito de términos. Al escribir
que la característica de Euler-Poincaré de este complejo es nula
\sref[0]{0.11.10.1}, se obtiene inmediatamente
\[
  \operatorname{EP}(\mathscr P_\bullet;y)
  =
  \operatorname{EP}(\mathscr P'_\bullet;y)
  +
  \operatorname{EP}(\mathscr P''_\bullet;y)
  \tag{7.9.6.1}\label{III.7.9.6.1}
\]
para todo $y\in Y$.

Supongamos, por ejemplo, que
$\mathscr P_\bullet=(\mathscr P_i)$ y que $\mathscr P_i=0$ para
$i<0$. Se tiene entonces la sucesión exacta de complejos
\[
\xymatrix@C=2.4em{
  \cdots \ar[r] & 0 \ar[r]\ar[d] & 0 \ar[r]\ar[d]
    & 0 \ar[r]\ar[d] & 0 \ar[r] & \cdots \\
  \cdots \ar[r] & 0 \ar[r]\ar[d] & 0 \ar[r]\ar[d]
    & \mathscr P_1 \ar[r]\ar[d] & \mathscr P_2 \ar[r] & \cdots \\
  \cdots \ar[r] & 0 \ar[r]\ar[d] & \mathscr P_0 \ar[r]\ar[d]
    & \mathscr P_1 \ar[r]\ar[d] & \mathscr P_2 \ar[r] & \cdots \\
  \cdots \ar[r] & 0 \ar[r]\ar[d] & \mathscr P_0 \ar[r]\ar[d]
    & 0 \ar[r]\ar[d] & 0 \ar[r] & \cdots \\
  \cdots \ar[r] & 0 \ar[r] & 0 \ar[r]
    & 0 \ar[r] & 0 \ar[r] & \cdots
}
\]
cuyas flechas verticales no nulas son los automorfismos identidad.
Aplicando \eqref{III.7.9.6.1} a esta sucesión exacta y razonando por
inducción sobre la longitud de $\mathscr P_\bullet$, se obtiene
\[
  \operatorname{EP}(\mathscr P_\bullet;y)
  =
  \sum_i(-1)^i\operatorname{EP}(\mathscr P_i;y).
  \tag{7.9.6.2}\label{III.7.9.6.2}
\]

\oldpage[III]{210}
Aquí, para todo $\mathscr O_X$-módulo coherente $\mathscr F$ plano
sobre $Y$, la notación $\operatorname{EP}(\mathscr F;y)$ (o
$\operatorname{EP}(f,\mathscr F;y)$, o
$\operatorname{EP}(Y,\mathscr F;y)$) significa
$\operatorname{EP}(\mathscr Q_\bullet;y)$ (con la notación relativa
correspondiente) para el complejo $\mathscr Q_\bullet$ cuyo único
término no nulo es $\mathscr F$ en grado $0$. Por tanto, el estudio
de las características de Euler-Poincaré de los complejos se reduce
al de los complejos con un solo término.
\end{env}

\begin{proposition}[7.9.7]
\label{III.7.9.7}
Bajo las hipótesis de \sref{III.7.9.4}, sean $Y'$ un preesquema
localmente noetheriano y $g:Y'\to Y$ un morfismo, pongamos
\[
  X'=X\times_Y Y',\qquad
  f'=f_{(Y')}:X'\longrightarrow Y',
\]
y sea
\[
  \mathscr P'_\bullet
  =
  \mathscr P_\bullet\otimes_{\mathscr O_Y}\mathscr O_{Y'}
\]
el complejo finito de $\mathscr O_{X'}$-módulos así obtenido.
Entonces $\mathscr P'_\bullet$ está formado por
$\mathscr O_{X'}$-módulos coherentes y planos sobre $Y'$, y, para
todo $y'\in Y'$, se tiene
\[
  \operatorname{EP}(\mathscr P'_\bullet;y')
  =
  \operatorname{EP}(\mathscr P_\bullet;g(y')).
  \tag{7.9.7.1}\label{III.7.9.7.1}
\]
\end{proposition}

\begin{proof}
Los $\mathscr O_{X'}$-módulos $\mathscr P'_i$, al ser imágenes
inversas de los $\mathscr P_i$ por la proyección $X'\to X$, son
coherentes. Son planos sobre $Y'$ en virtud de
\sref[0]{0.6.2.1} y \sref[I]{I.4.14.5}, pues la cuestión es local
en $X$, $Y$ e $Y'$. Por último, $f'$ es propio por
\sref[II]{II.5.4.2}, de modo que el primer miembro de
\eqref{III.7.9.7.1} está definido. La fórmula
\eqref{III.7.9.7.1} se deduce entonces de
\textup{(6.10.4.2)}, \sref{III.7.7.2} y
\sref{III.7.6.7}, reduciéndose, como siempre se puede hacer, al
caso en que $Y$ e $Y'$ son afines.
\end{proof}

\begin{proposition}[7.9.8]
\label{III.7.9.8}
Supongamos que se cumplen las hipótesis de \sref{III.7.9.4} y,
además, que existe un entero $i_0$ tal que
\[
  \mathscr T_i(\kres(y))=0
  \qquad\text{para $i\neq i_0$ y todo $y\in Y$.}
\]
Entonces
\[
  \mathscr T_{i_0}(\mathscr O_Y)
  =
  \mathscr H^{-i_0}(f,\mathscr P^\bullet)
\]
es un $\mathscr O_Y$-módulo localmente libre, cuyo rango en
$y\in Y$ es
\[
  (-1)^{i_0}\operatorname{EP}(f,\mathscr P_\bullet;y).
\]
\end{proposition}

\begin{proof}
Observemos primero que los funtores
$\mathscr T_i^{\mathscr O_y}$ cumplen las hipótesis de
\sref{III.7.4.4}, por lo que \sref{III.7.4.7} les es aplicable.
La hipótesis y \sref{III.7.5.3} implican que
$\mathscr T_i^{\mathscr O_y}$ es nulo para $i\neq i_0$.
En consecuencia, $\mathscr T_{i_0}$ es exacto por
\sref{III.7.3.3}, y \sref{III.7.8.4} muestra que
$\mathscr H^{-i_0}(f,\mathscr P^\bullet)$ es localmente libre.
Su rango en $y\in Y$ es
\[
  \operatorname{rank}_{\kres(y)}
    \mathscr T_{i_0}(\kres(y))
  =
  (-1)^{i_0}\operatorname{EP}(f,\mathscr P_\bullet;y),
\]
por definición, ya que $\mathscr T_i(\kres(y))=0$ para
$i\neq i_0$.
\end{proof}

\begin{corollary}[7.9.9]
\label{III.7.9.9}
Sean $Y$ un preesquema localmente noetheriano,
$f:X\to Y$ un morfismo propio y $\mathscr F$ un
$\mathscr O_X$-módulo coherente y plano sobre $Y$. Supongamos que
existe un entero $i_0$ tal que
\[
  H^i\bigl(f^{-1}(y),
    \mathscr F\otimes_{\mathscr O_Y}\kres(y)\bigr)=0
\]
para todo $i\neq i_0$ y todo $y\in Y$. Entonces
$R^{i_0}f_*(\mathscr F)$ es un $\mathscr O_Y$-módulo localmente
libre, cuyo rango en $y$ es
$(-1)^{i_0}\operatorname{EP}(f,\mathscr F;y)$.
\end{corollary}

En particular:

\begin{corollary}[7.9.10]
\label{III.7.9.10}
Bajo las hipótesis preliminares de \sref{III.7.9.9} relativas a
$X$, $Y$ y $\mathscr F$, supongamos que
$R^if_*(\mathscr F)=0$ para todo $i>0$. Entonces
$f_*(\mathscr F)$ es un $\mathscr O_Y$-módulo localmente libre,
cuyo rango en $y$ es $\operatorname{EP}(f,\mathscr F;y)$.
\end{corollary}

\begin{proof}
En virtud de \sref{III.7.9.9}, basta demostrar el lema siguiente.
\end{proof}

\begin{lemma}[7.9.10.1]
\label{III.7.9.10.1}
Bajo las hipótesis de \sref{III.7.9.10}, se tiene
\[
  H^i\bigl(f^{-1}(y),
    \mathscr F\otimes_{\mathscr O_Y}\kres(y)\bigr)=0
\]
para todo $i>0$ y todo $y\in Y$.
\end{lemma}

\begin{proof}
Podemos limitarnos al caso en que $Y=\Spec(A)$ es afín.
Con las notaciones de \sref{III.7.9.4}, y estando
$\mathscr P_\bullet$ reducido a su único término no nulo
$\mathscr F$ en grado $0$, la hipótesis da
$\mathscr T_p(\mathscr O_Y)=0$ para $p<0$. Se deduce de
\sref{III.7.3.7} que $\mathscr T_p$ es exacto para $p<0$, y el
lema resulta entonces de la equivalencia de las condiciones a) y
d) de \sref{III.7.7.5}.
\end{proof}

\oldpage[III]{211}
\begin{proposition}[7.9.11]
\label{III.7.9.11}
Bajo las hipótesis de \sref{III.7.9.4}, sea $\mathscr L$ un
$\mathscr O_X$-módulo inversible muy amplio sobre $Y$ y, para todo
$n\in\bb Z$, pongamos
\[
  \mathscr P_\bullet(n)
  =
  \mathscr P_\bullet
    \otimes_{\mathscr O_X}\mathscr L^{\otimes n}.
\]
Entonces, para todo $y\in Y$, la función
\[
  n\longmapsto
  \operatorname{EP}(f,\mathscr P_\bullet(n);y)
  \tag{7.9.11.1}\label{III.7.9.11.1}
\]
es un polinomio con coeficientes racionales, y es el mismo para
todos los puntos de una misma componente conexa de $Y$.
\end{proposition}

\begin{proof}
Es claro que $\mathscr P_\bullet(n)$ es un complejo de
$\mathscr O_X$-módulos planos sobre $Y$. En virtud de
\eqref{III.7.9.6.2}, podemos reducirnos al caso en que
$\mathscr P_\bullet$ tiene un único término no nulo
$\mathscr F\neq0$ en grado $0$. Como se trata de cuestiones
locales sobre $Y$, podemos suponer además que $Y$ es afín y que
$f$ es proyectivo \sref[II]{II.5.5.3}. Pongamos
\[
  X_y=f^{-1}(y),\qquad
  \mathscr L_y=\mathscr L\otimes_{\mathscr O_Y}\kres(y);
\]
entonces $\mathscr L_y$ es un $\mathscr O_{X_y}$-módulo muy amplio
\sref[II]{II.4.4.10}. Por \sref{III.7.7.2}, para el funtor
$\mathscr T_\bullet$ asociado a $\mathscr P_\bullet(n)$ se tiene
\[
  \mathscr T_i(\kres(y))
  =
  H^{-i}\bigl(
    X_y,\mathscr F_y\otimes\mathscr L_y^{\otimes n}\bigr),
  \qquad
  \mathscr F_y=\mathscr F\otimes_{\mathscr O_Y}\kres(y).
\]
Por consiguiente, $\operatorname{EP}(f,\mathscr F(n);y)$ no es
sino la característica de Euler-Poincaré
$\chi_{\kres(y)}(\mathscr F_y(n))$ definida en
\sref{III.2.5.1}. De \sref{III.2.5.3} se deduce que
\eqref{III.7.9.11.1} es un polinomio. Además, para cada $n$ fijo,
su valor es constante en cada componente conexa de $Y$ por
\sref{III.7.9.4}; esto termina la demostración.
\end{proof}

\begin{env}[7.9.12]
\label{III.7.9.12}
El polinomio \eqref{III.7.9.11.1}, con coeficientes racionales, se
denotará por
\[
  \operatorname{PH}(f,\mathscr P_\bullet;y)
  \qquad\text{o}\qquad
  \operatorname{PH}(\mathscr P_\bullet;y),
\]
y se llamará \emph{polinomio de Hilbert en $y$} relativo a
$\mathscr P_\bullet$, $f$ y $\mathscr L$ (o simplemente
\emph{polinomio de Hilbert en $y$} de $\mathscr P_\bullet$, o de
$f$, cuando no haya riesgo de confusión). Cuando $Y$ es conexo y
no vacío, se suprime la referencia a $y$ en la notación y la
terminología. Este invariante desempeñará un papel esencial en el
capítulo~V, en la teoría de los «móduli» de los haces cociente
coherentes de un haz coherente dado.

Con las notaciones de \sref{III.7.9.6} y
\sref{III.7.9.11}, se tiene
\[
  \operatorname{PH}(\mathscr P_\bullet;y)
  =
  \operatorname{PH}(\mathscr P'_\bullet;y)
  +
  \operatorname{PH}(\mathscr P''_\bullet;y),
  \tag{7.9.12.1}\label{III.7.9.12.1}
\]
y, en particular,
\[
  \operatorname{PH}(\mathscr P_\bullet;y)
  =
  \sum_i(-1)^i\operatorname{PH}(\mathscr P_i;y).
  \tag{7.9.12.2}\label{III.7.9.12.2}
\]
Estas fórmulas se deducen inmediatamente de
\eqref{III.7.9.6.1} y \eqref{III.7.9.6.2}. Asimismo, con las
notaciones e hipótesis de \sref{III.7.9.7}, se tiene
\[
  \operatorname{PH}(\mathscr P'_\bullet;y')
  =
  \operatorname{PH}(\mathscr P_\bullet;g(y')).
  \tag{7.9.12.3}\label{III.7.9.12.3}
\]
\end{env}

La fórmula \eqref{III.7.9.12.2} reduce el estudio de los
polinomios de Hilbert de un complejo al de los polinomios de
Hilbert de un solo $\mathscr O_X$-módulo plano sobre $Y$.
Estos últimos admiten la siguiente interpretación notable,
independiente de consideraciones homológicas.

\begin{corollary}[7.9.13]
\label{III.7.9.13}
Sean $Y$ un preesquema noetheriano, $f:X\to Y$ un morfismo propio,
$\mathscr L$ un $\mathscr O_X$-módulo inversible muy amplio sobre
$Y$ y $\mathscr F$ un $\mathscr O_X$-módulo coherente y plano
sobre $Y$. Existe un entero $n_0$ tal que, para $n\geq n_0$,
$f_*(\mathscr F(n))$ es un $\mathscr O_Y$-módulo localmente libre
cuyo rango en $y\in Y$ es
\[
  \operatorname{PH}(f,\mathscr F;y)(n).
\]
\end{corollary}

\begin{proof}
Como $f$ es proyectivo \sref[II]{II.5.5.3}, existe un entero $n_0$
tal que, para $n\geq n_0$,

\oldpage[III]{212}
\[
  R^if_*(\mathscr F(n))=0
  \qquad\text{para todo $i>0$}
\]
por \sref{III.2.2.1}. La conclusión se deduce entonces de
\sref{III.7.9.10}.
\end{proof}

El criterio de planitud siguiente será importante en el
capítulo~V, en la teoría de los «móduli» de los haces coherentes.

\begin{proposition}[7.9.14]
\label{III.7.9.14}
Sean $Y$ un preesquema noetheriano, $f:X\to Y$ un morfismo
proyectivo y $\mathscr L$ un $\mathscr O_X$-módulo inversible
amplio para $f$. Para todo $\mathscr O_X$-módulo $\mathscr F$ y
todo $n\in\bb Z$, pongamos
\[
  \mathscr F(n)=\mathscr F\otimes\mathscr L^{\otimes n}.
\]
Para que un $\mathscr O_X$-módulo coherente $\mathscr F$ sea plano
sobre $Y$, es necesario y suficiente que exista un entero $n_0$
tal que, para todo $n\geq n_0$, $f_*(\mathscr F(n))$ sea un
$\mathscr O_Y$-módulo localmente libre.
\end{proposition}

\begin{proof}
La necesidad de la condición se demuestra como en
\sref{III.7.9.13}; el resultado de \sref{III.2.2.1} se aplica a un
haz amplio porque $f$ es proyectivo.

Recíprocamente, podemos reducirnos al caso en que
$Y=\Spec(A)$ es afín. En virtud de la hipótesis y de
\sref{III.2.2.2}{(i)}, los $A$-módulos
$\Gamma(X,\mathscr F(n))$ son proyectivos y de tipo finito
(Bourbaki, \emph{Alg.~comm.}, cap.~II, \S5,
n.\textsuperscript{o}~2, teorema~1). Sea
\[
  S=\bigoplus_{n\geq0}\Gamma(X,\mathscr L^{\otimes n})
\]
el anillo graduado asociado. Se sabe que $X$ se identifica
canónicamente con $\Proj(S)$ por
\sref[II]{II.4.5.2}{(b)} y \sref[II]{II.5.4.4}. Sea
\[
  M=\bigoplus_{n\geq n_0}\Gamma(X,\mathscr F(n)).
\]
Sustituyendo, si es preciso, $\mathscr L$ por una potencia
$\mathscr L^{\otimes d}$, podemos suponer que $S$ está generado por
un número finito de elementos de grado $1$
\sref{III.2.3.5.1}; se deduce entonces de
\sref[II]{II.2.7.5} y \sref[II]{II.2.7.2} que $\mathscr F$ se
identifica con $\mathscr P\!roj_0(M)$.

Para todo elemento homogéneo $g\in S$ de grado positivo, se tiene,
por tanto,
\[
  \Gamma(X_g,\mathscr F)=M_{(g)}.
\]
Ahora bien, $M$, al ser una suma directa de $A$-módulos
proyectivos, es un $A$-módulo plano. Por consiguiente, $M_g$ es
plano por \sref[0]{0.6.3.2}, y también lo es $M_{(g)}$, que es un
factor directo de $M_g$ por \sref[0]{0.6.1.2}. Se concluye de
\sref[I]{I.4.14.5} que $\mathscr F$ es plano sobre $Y$ en todo
punto de $X_g$. Como los abiertos $X_g$ recubren $X$, queda
demostrada la proposición.
\end{proof}

\bigskip
\noindent\emph{(Continuará.)}

\clearpage
\oldpage[III]{213}
